\chapter{Vektorer og matricer}\label{chap:ch8}
En \textbf{vektor} er et matematisk objekt, der har både en \textbf{størrelse} (længde) og en \textbf{retning}. I to eller tre dimensioner tegnes en vektor ofte som en pil fra ét punkt til et andet. For eksempel beskriver vektoren \(\vect{v} = (x_1, x_2)\) i 2D-rummet en bevægelse fra origo \((0,0)\) til punktet \((x_1, x_2)\).

Vektorer bruges til at beskrive størrelser som hastighed, kraft og forskydning, der kræver både en størrelse og en retning for at være fuldt beskrevet. Vektorer kan lægges sammen, skaleres med et tal og opløses i komposanter.

Matricer er derimod rektangulære skemaer af tal ordnet i rækker og søjler. De bruges til at beskrive og løse lineære ligningssystemer, udføre transformationer og beskrive sammenhænge mellem mængder af vektorer.

\section{Vektorer i \texorpdfstring{$\mathbb{R}^2$}{R2} og \texorpdfstring{$\mathbb{R}^n$}{Rn}}

I lineær algebra skrives vektorer ofte som \textbf{søjlematricer}. For eksempel kan vektoren \(\vect{v} = (x_1, x_2)\) skrives som en søjlematrix:

\[ \vect{v} = \begin{bmatrix} x_1 \\ x_2 \end{bmatrix} \]

Vi siger, at to vektorer er \textbf{ens}, hvis og kun hvis deres tilsvarende indgange er ens.

\begin{example}
    Følgende er vektorer i $\mathbb{R}^2$ (dvs. planen bestående af ordnede par af reelle tal):

    \[
    \vect{u}=\left[\begin{array}{l}
    3 \\
    5
    \end{array}\right], \quad \vect{v}=\left[\begin{array}{l}
    5 \\
    3
    \end{array}\right]
    \]
    De er ikke ens, fordi deres tilsvarende indgange ikke stemmer overens.
\end{example}

\begin{figure}[htbp]
    \centering
    % First graph (points)
    \begin{minipage}{0.45\textwidth}
        \centering
        \begin{tikzpicture}
            % Draw the axes
            \draw[->] (-3, 0) -- (3, 0) node[right] {$x_1$};
            \draw[->] (0, -3) -- (0, 3) node[above] {$x_2$};

            % Plot the points
            \filldraw[headercolor] (-2,-1) circle (2pt) node[anchor=north east] {(-2, -1)};
            \filldraw[headercolor] (1,1) circle (2pt) node[anchor=south west] {(1, 1)};
            \filldraw[headercolor] (2,-2) circle (2pt) node[anchor=north west] {(2, -2)};
        \end{tikzpicture}
        \subcaption{Punkter i planen}
    \end{minipage}
    \hfill
    % Second graph (vectors)
    \begin{minipage}{0.45\textwidth}
        \centering
        \begin{tikzpicture}
            % Draw the axes
            \draw[->] (-3, 0) -- (3, 0) node[right] {$x_1$};
            \draw[->] (0, -3) -- (0, 3) node[above] {$x_2$};

            % Plot the vectors
            \draw[->, thick, headercolor] (0,0) -- (-2,-1) node[anchor=north east] {(-2, -1)};
            \draw[->, thick, headercolor] (0,0) -- (1,1) node[anchor=south west] {(1, 1)};
            \draw[->, thick, headercolor] (0,0) -- (2,-2) node[anchor=north west] {(2, -2)};
        \end{tikzpicture}
        \subcaption{Vektorer ud fra origo}
    \end{minipage}
    \caption{Punkter og vektorer i planen}
    \label{fig:points-vectors}
\end{figure}

Summen af to vektorer $\vect{u}$ og $\vect{v}$ i $\mathbb{R}^2$, som skrives $\vect{u} + \vect{v}$, fås ved at lægge de tilsvarende indgange i $\vect{u}$ og $\vect{v}$ sammen. Givet et reelt tal $c$ fås multiplummet af $\vect{u}$ med skalaren $c$, som skrives $c \vect{u}$, ved at gange hver indgang i $\vect{u}$ med $c$.

\begin{example}
    Hvis $\vect{u}$ og $\vect{v}$ er som i det foregående eksempel, så er 

$$
\begin{aligned}
\vect{u} + \vect{v} & =\left[\begin{array}{l}
3 \\
5
\end{array}\right] + \left[\begin{array}{l}
5 \\
3
\end{array}\right] = \left[\begin{array}{l}
3+5 \\
5+3
\end{array}\right] = \left[\begin{array}{l}
8 \\
8
\end{array}\right], \text { og } \\
6 \vect{u} & = 6\left[\begin{array}{l}
3 \\
5
\end{array}\right] = \left[\begin{array}{l}
6 \cdot 3 \\
6 \cdot 5
\end{array}\right] = \left[\begin{array}{l}
18 \\
30
\end{array}\right]
\end{aligned}
$$

\end{example}

\begin{remark}
    Det er ofte en hjælp at identificere en vektor $\left[\begin{array}{l}a \\ b\end{array}\right] \in \mathbb{R}^2$ med et geometrisk punkt $(a, b)$ i planen for at få et billede af det, vi arbejder med. Se \autoref{fig:points-vectors} for et eksempel.
\end{remark}

Hvis vi opfatter $\vect{u}$ og $\vect{v}$ i $\mathbb{R}^2$ som punkter i planen, svarer $\vect{u} + \vect{v}$ til det fjerde hjørne i det parallelogram, hvis øvrige hjørner er $\vect{0}, \vect{u}$ og $\vect{v}$.
Bemærk: Med $\vect{0}$ mener vi nulvektoren, dvs. vektoren, hvis indgange alle er nul. I $\mathbb{R}^2$ har vi $\vect{0} = \left[\begin{array}{l}
0 \\
0
\end{array}\right]$.

Vi kalder dette \textit{parallelogramreglen for addition}, og den er vist i \autoref{fig:parallelogram-law}.

\begin{figure}[htbp]
    \centering
    \begin{tikzpicture}
        % Draw axes
        \draw[->] (-0.5, 0) -- (4, 0) node[right] {$x_1$};
        \draw[->] (0, -0.5) -- (0, 3) node[above] {$x_2$};

        % Define the vectors u and v
        \coordinate (O) at (0, 0);
        \coordinate (U) at (1.5, 2);
        \coordinate (V) at (3, 0.5);
        \coordinate (UplusV) at (3+1.5, 2+0.5);

        % Draw vectors
        \draw[->, thick] (O) -- (U);
        \draw[->, thick] (O) -- (V);
        \draw[->, thick, headercolor] (O) -- (UplusV) node[midway, right] {$\vect{u} + \vect{v}$};

        % Draw the parallelogram
        \draw[dashed] (U) -- (UplusV);
        \draw[dashed] (V) -- (UplusV);
        \fill[headercolor, opacity=0.3] (O) -- (U) -- (UplusV) -- (V) -- cycle;

        % Add labels
        \node[above right] at (UplusV) {$\vect{u} + \vect{v}$};
        \node[above] at (U) {$\vect{u}$};
        \node[below right] at (V) {$\vect{v}$};

        % Mark the origin
        \node[below left] at (O) {$\vect{0}$};
    \end{tikzpicture}
    \caption{Parallelogramreglen for vektoraddition i $\mathbb{R}^2$.}
    \label{fig:parallelogram-law}
\end{figure}

Disse idéer kan generaliseres til rum af højere dimension. Mere præcist kan vi definere $\mathbb{R}^n$ på følgende måde.

\begin{definition}
    For hvert positivt heltal $n$ lader vi $\mathbb{R}^n$ betegne samlingen af ordnede $n$-tupler, hvor hver indgang ligger i $\mathbb{R}$. Vi skriver ofte disse elementer som $n \times 1$-matricer. Addition og multiplikation med skalar af vektorer i $\mathbb{R}^n$ defineres på samme måde som for $\mathbb{R}^2$. Det vil sige, at vi regner koordinat for koordinat.
\end{definition}

\begin{example}
    Hvis $u_1, u_2, \ldots, u_n \in \mathbb{R}$, så er

    \[
    \vect{u} = \left[\begin{array}{c}
    u_1 \\
    u_2 \\
    \vdots \\
    u_n
    \end{array}\right] \in \mathbb{R}^n.
    \]
     
\end{example}

\begin{example}
    Hvis $\vect{u}$ og $\vect{v}$ ligger i $\mathbb{R}^n$ (med indgange betegnet henholdsvis $u_1, \ldots, u_n$ og $v_1, \ldots, v_n$), og $c \in \mathbb{R}$, så er

    \[
    \begin{aligned}
    & \vect{u} + \vect{v} = {\left[\begin{array}{c}
    u_1 \\
    u_2 \\
    \vdots \\
    u_n
    \end{array}\right] + \left[\begin{array}{c}
    v_1 \\
    v_2 \\
    \vdots \\
    v_n
    \end{array}\right] = \left[\begin{array}{c}
    u_1 + v_1 \\
    u_2 + v_2 \\
    \vdots \\
    u_n + v_n
    \end{array}\right], \text { og } } 
    & c \vect{u} = c\left[\begin{array}{c}
    u_1 \\
    u_2 \\
    \vdots \\
    u_n
    \end{array}\right] = \left[\begin{array}{c}
    c u_1 \\
    c u_2 \\
    \vdots \\
    c u_n
    \end{array}\right]
    \end{aligned}
    \]   
\end{example}

\begin{custombox}{Algebraiske egenskaber ved $\mathbb{R}^n$}
    Lad $\vect{u}, \vect{v}, \vect{w} \in \mathbb{R}^n$ og $c, d \in \mathbb{R}$. Så gælder følgende egenskaber:
    \begin{enumerate}
        \item \textbf{Den kommutative lov for addition:} $\vect{u} + \vect{v} = \vect{v} + \vect{u}$.
        \item \textbf{Den associative lov for addition:} $(\vect{u} + \vect{v}) + \vect{w} = \vect{u} + (\vect{v} + \vect{w})$.
        \item \textbf{Neutralt element for addition:} Der findes en vektor $\vect{0} \in \mathbb{R}^n$, så $\vect{u} + \vect{0} = \vect{u}$ for alle $\vect{u} \in \mathbb{R}^n$.
        \item \textbf{Modsat vektor:} For hver $\vect{u} \in \mathbb{R}^n$ findes en vektor $-\vect{u} \in \mathbb{R}^n$, så $\vect{u} + (-\vect{u}) = \vect{0}$.
        \item \textbf{Den distributive lov:} $c(\vect{u} + \vect{v}) = c\vect{u} + c\vect{v}$ og $(c + d)\vect{u} = c\vect{u} + d\vect{u}$.
        \item \textbf{Den associative lov for multiplikation med skalar:} $c(d\vect{u}) = (cd)\vect{u}$.
        \item \textbf{Neutralt element for multiplikation:} $1\vect{u} = \vect{u}$.
    \end{enumerate}  
\end{custombox}

En \textbf{linearkombination} er en måde at kombinere vektorer på ved hjælp af multiplikation med skalar og addition. Givet en mængde af vektorer ganger vi hver vektor med en skalar og lægger derefter resultaterne sammen. Linearkombinationer hjælper os med at forstå, hvordan vektorer hænger sammen, og om en vektor kan udtrykkes ved hjælp af andre.

\begin{definition}{Linearkombinationer}
    Givet en mængde af vektorer $\vect{v}_1, \vect{v}_2, \ldots, \vect{v}_p \in \mathbb{R}^n$ og skalarer $c_1, c_2, \ldots, c_p \in \mathbb{R}$ kaldes vektoren $\vect{y}$ givet ved

\[
\vect{y} = c_1 \vect{v}_1 + \cdots + c_p \vect{v}_p
\]

en linearkombination af $\vect{v}_1, \vect{v}_2, \ldots, \vect{v}_p$ med vægtene $c_1, c_2, \ldots, c_p$.

\end{definition}

\begin{example}
    Lad $\vect{v}_1 = \left[\begin{array}{l}
        1 \\
        2
        \end{array}\right]$ og $\vect{v}_2 = \left[\begin{array}{c}
        1 \\
        -1
        \end{array}\right]$. Nogle linearkombinationer af $\vect{v}_1$ og $\vect{v}_2$ er
        
        \[
        \begin{aligned}
        \vect{0} & = 0 \vect{v}_1 + 0 \vect{v}_2 \\
        {\left[\begin{array}{l}
        3 \\
        0
        \end{array}\right] } & = \vect{v}_1 + 2 \vect{v}_2 \\
        {\left[\begin{array}{l}
        -5 \\
        -1
        \end{array}\right] } & = -2 \vect{v}_1 - 3 \vect{v}_2, \text { og } \\
        {\left[\begin{array}{l}
        5 \\
        1
        \end{array}\right] } & = 2 \vect{v}_1 + 3 \vect{v}_2
        \end{aligned}
        \]
        
\end{example}

\subsection*{Vektorligninger og lineære ligningssystemer}

Vi ønsker ofte at vide, om en vektor $\vect{b}$ kan dannes som en linearkombination af en anden mængde af vektorer $\vect{a}_1, \ldots, \vect{a}_n$. Fremgangsmåden til at afgøre det er givet ved følgende.

\begin{custombox}{Brug af matricer til at bestemme linearkombinationer}
En vektorligning

\[
x_1 \vect{a}_1 + x_2 \vect{a}_2 + \cdots + x_n \vect{a}_n = \vect{b}
\]

har samme løsningsmængde som det lineære ligningssystem, hvis totalmatrix er

\[
\left[\begin{array}{lllll}
\vect{a}_1 & \vect{a}_2 & \ldots & \vect{a}_n & \vect{b}
\end{array}\right].
\]    
\end{custombox}

Mere præcist: Hvis vektorerne $\vect{a}_i$ ligger i $\mathbb{R}^m$ med

\[
\vect{a}_1 = \left[\begin{array}{c}
a_{11} \\
a_{21} \\
\vdots \\
a_{m1}
\end{array}\right], \quad \vect{a}_2 = \left[\begin{array}{c}
a_{12} \\
a_{22} \\
\vdots \\
a_{m2}
\end{array}\right], \ldots, \quad \vect{a}_n = \left[\begin{array}{c}
a_{1n} \\
a_{2n} \\
\vdots \\
a_{mn}
\end{array}\right], \quad \text{og} \quad \vect{b} = \left[\begin{array}{c}
b_1 \\
b_2 \\
\vdots \\
b_m
\end{array}\right]
\]

så rækkereducerer man matricen

\[
\begin{aligned}
& \left[\begin{array}{ccccc}
a_{11} & a_{12} & \ldots & a_{1n} & b_1 \\
a_{21} & a_{22} & \ldots & a_{2n} & b_2 \\
\vdots & \vdots & \ddots & \vdots & \vdots \\
a_{m1} & a_{m2} & \ldots & a_{mn} & b_m
\end{array}\right] \\
& \begin{array}{ccccc}
\hspace{0.5cm} \uparrow & \hspace{0.3cm} \uparrow & \hspace{0.0cm} \ldots & \hspace{0.2cm} \uparrow & \hspace{0.2cm} \uparrow \\
\hspace{0.4cm} \vect{a}_1 & \hspace{0.3cm} \vect{a}_2 & \hspace{0.1cm} \ldots & \hspace{0.2cm} \vect{a}_n & \hspace{0.2cm} \vect{b}
\end{array}
\end{aligned}
\]

for at afgøre, om der findes et sæt vægte $x_1, \ldots, x_n$, der passer.

\begin{example} Lad
\[
\begin{aligned}
&\vect{a}_1=\left[\begin{array}{c}
2 \\
-1 \\
1
\end{array}\right], \quad \vect{a}_2=\left[\begin{array}{c}
0 \\
8 \\
-2
\end{array}\right], \quad \vect{a}_3=\left[\begin{array}{l}
6 \\
5 \\
1
\end{array}\right], \quad \text { og } \vect{b}=\left[\begin{array}{c}
10 \\
3 \\
7
\end{array}\right]
\end{aligned}
\]

Vi afgør, om $\vect{b}$ er en linearkombination af $\vect{a}_1, \vect{a}_2, \vect{a}_3$, dvs. om der findes et sæt vægte $x_1, x_2, x_3$, så $x_1 \vect{a}_1 + x_2 \vect{a}_2 + x_3 \vect{a}_3 = \vect{b}$. Ifølge ovenstående oversætter vi spørgsmålet til matrixform.

\[
\begin{alignedat}{3}
& \left[\begin{array}{cccc}
2 & 0 & 6 & 10 \\
-1 & 8 & 5 & 3 \\
1 & -2 & 1 & 7
\end{array}\right]
& \quad & r_1 \leftrightarrow r_3 
& \quad & 
\left[\begin{array}{cccc}
1 & -2 & 1 & 7 \\
-1 & 8 & 5 & 3 \\
2 & 0 & 6 & 10
\end{array}\right] \\[10pt]
& & \quad & \begin{aligned}
    r_2 & \mapsto r_2 + r_1 \\
    r_3 & \mapsto r_3 - 2r_1
\end{aligned}
& \quad & 
\left[\begin{array}{cccc}
1 & -2 & 1 & 7 \\
0 & 6 & 6 & 10 \\
0 & 4 & 4 & -4
\end{array}\right] \\[10pt]
& & \quad & r_3 \mapsto r_3 - \frac{4}{6} r_2
& \quad & 
\left[\begin{array}{cccc}
1 & -2 & 1 & 7 \\
0 & 6 & 6 & 10 \\
0 & 0 & 0 & \blacksquare
\end{array}\right] .
\end{alignedat}
\]

Da symbolet $\blacksquare$ betegner noget, der er forskelligt fra nul, ser vi, at der ikke er nogen løsning, dvs. at $\vect{b}$ ikke er en linearkombination af $\vect{a}_1, \vect{a}_2, \vect{a}_3$.

\end{example}

\begin{example} Lad

\[
\vect{a}_1 = \left[\begin{array}{l}
1 \\
0 \\
1
\end{array}\right], \quad
\vect{a}_2 = \left[\begin{array}{c}
-4 \\
6 \\
-4
\end{array}\right], \quad
\vect{a}_3 = \left[\begin{array}{c}
-6 \\
7 \\
5
\end{array}\right], \quad
\text{og } \vect{b} = \left[\begin{array}{c}
11 \\
-5 \\
9
\end{array}\right]
\]

Vi afgør, om $\vect{b}$ er en linearkombination af $\vect{a}_1, \vect{a}_2, \vect{a}_3$. Vi reducerer den tilhørende matrix til echelonform:

\[
\left[\begin{array}{cccc}
1 & -4 & -6 & 11 \\
0 & 6 & 7 & -5 \\
1 & -4 & 5 & 9
\end{array}\right]
\quad r_3 \mapsto r_3 - r_1 \quad
\left[\begin{array}{cccc}
1 & -4 & -6 & 11 \\
0 & 6 & 7 & -5 \\
0 & 0 & 11 & -2
\end{array}\right]
\]

og vi ser, at der findes en løsning. Nu finder vi de vægte $x_1, x_2, x_3$, der passer, ved at bestemme den reducerede echelonform.

\[
\begin{alignedat}{3}
& r_3 \mapsto \frac{1}{11} r_3 \quad &
& \left[\begin{array}{cccc}
1 & -4 & -6 & 11 \\
0 & 6 & 7 & -5 \\
0 & 0 & 1 & -\frac{2}{11}
\end{array}\right] \\[10pt]
& \begin{aligned}
    r_2 & \mapsto r_2 - 7r_3 \\
    r_1 & \mapsto r_1 + 6r_3
\end{aligned} & \quad &
\left[\begin{array}{cccc}
1 & -4 & 0 & \frac{109}{11} \\
0 & 6 & 0 & -\frac{41}{11} \\
0 & 0 & 1 & -\frac{2}{11}
\end{array}\right]
\end{alignedat}
\]

\[
\begin{alignedat}{3}
& r_2 \mapsto \frac{1}{6}r_2 \quad &
& \left[\begin{array}{cccc}
1 & -4 & 0 & \frac{109}{11} \\
0 & 1 & 0 & -\frac{41}{66} \\
0 & 0 & 1 & -\frac{2}{11}
\end{array}\right] \\[10pt]
& r_1 \mapsto r_1 + 4r_2 \quad &
& \left[\begin{array}{cccc}
1 & 0 & 0 & \frac{245}{33} \\
0 & 1 & 0 & -\frac{41}{66} \\
0 & 0 & 1 & -\frac{2}{11}
\end{array}\right],
\end{alignedat}
\]

så $\left(x_1, x_2, x_3\right)=\left(\frac{245}{33},-\frac{41}{66},-\frac{2}{11}\right)$. Da der ikke er nogen frie variable, er dette den entydige løsning.

\end{example}

\textbf{Udspændet} (Span) af en mængde af vektorer er samlingen af alle mulige linearkombinationer af vektorerne. Det er med andre ord mængden af alle de vektorer, man kan nå ved at skalere og lægge de givne vektorer sammen (se eksempler i \autoref{fig:span_geometry}). Udspændet giver os indsigt i, hvilket »rum« vektorerne dækker:

\begin{definition}{Udspændet af en mængde af vektorer}
    Hvis $\vect{v}_1, \ldots, \vect{v}_p$ ligger i $\mathbb{R}^n$, betegnes mængden af alle linearkombinationer af $\vect{v}_1, \ldots, \vect{v}_p$ med \text{Span} $\left\{\vect{v}_1, \ldots, \vect{v}_p\right\}$, og den kaldes den delmængde af $\mathbb{R}^n$, der udspændes af $\vect{v}_1, \ldots, \vect{v}_p$. Udspændet af $\vect{v}_1, \ldots, \vect{v}_p$ består med andre ord af alle vektorer, der kan skrives på formen
    $$
    c_1 \vect{v}_1 + \cdots + c_p \vect{v}_p
    $$
    hvor $c_1, \ldots, c_p$ er skalarer.

\end{definition}

\begin{figure}[htbp]
    \centering

    \begin{minipage}{0.45\textwidth}
        \centering
        \begin{tikzpicture}[scale = 0.4]
            \useasboundingbox (-8,-8) rectangle (8,8);

            \def\A{7}
            \draw[->] (-\A,0) -- (\A,0) node[right] {$x_1$};
            \draw[->] (0,-\A) -- (0,\A) node[above] {$x_2$};

            \def\ux{2}
            \def\uy{1}

            \draw[mseViaBlue] (-\A,-\uy/\ux*\A) -- (\A,\uy/\ux*\A);

            \foreach \s in {-3,-2,-1,1,2,3} {
                \fill[mseViaBlue] ( {\s*\ux}, {\s*\uy} ) circle (3pt)
                node[mseViaBlue, anchor=north, xshift=2mm] {$\s\vect u$};
            }

            \draw[->, very thick] (0,0) -- (\ux,\uy)
            node[midway, yshift=2mm] {$\vect u$};

            \node at (4,5) {$\operatorname{span}\{\vect u\}$};
        \end{tikzpicture}

        \subcaption{Udspændet af én vektor i $\mathbb{R}^2$ er en linje gennem origo.}
    \end{minipage}
    \hfill
    \begin{minipage}{0.45\textwidth}
        \centering

        \tdplotsetmaincoords{60}{120}

        \begin{tikzpicture}[tdplot_main_coords, scale=0.4]
            \begin{scope}[tdplot_screen_coords]
                \path[use as bounding box] (-8,-8) rectangle (8,8);
            \end{scope}

            \coordinate (U) at (4.4,2.0,0.8);
            \coordinate (V) at (1.2,-2.8,2.8);

            \def\A{8}
            \draw[->, gray!50] (0,0,0) -- (\A,0,0)  node[ below] {$x_1$};
            \draw[gray!50]     (0,0,0) -- (0,-\A,0);
            \draw[gray!50]     (0,0,0) -- (0,0,-\A);

            \def\L{1.8}

            \coordinate (P1) at ($ -\L*(U) -\L*(V) $);
            \coordinate (P2) at ($  \L*(U) -\L*(V) $);
            \coordinate (P3) at ($  \L*(U) +\L*(V) $);
            \coordinate (P4) at ($ -\L*(U) +\L*(V) $);

            \fill[mseViaBlue, opacity=0.30] (P1) -- (P2) -- (P3) -- (P4) -- cycle;

            \draw (0,0,0) -- (-\A,0,0);
            \draw[->] (0,0,0) -- (0,\A,0) node[right] {$x_2$};
            \draw[->] (0,0,0) -- (0,0,\A) node[above] {$x_3$};

            \begin{scope}
                \clip (P1) -- (P2) -- (P3) -- (P4) -- cycle;

                \foreach \k in {-4,-3,...,4} {
                    \draw[mseViaBlue!70, dashed]
                    ($ -\L*(U) + \k*(V) $) -- ($ \L*(U) + \k*(V) $);

                    \draw[mseViaBlue!70, dashed]
                    ($ \k*(U) -\L*(V) $) -- ($ \k*(U) + \L*(V) $);
                }
            \end{scope}

            \foreach \a/\b in {1/1, 1/-1, 0/-1} {
                \fill[mseViaBlue] ($ \a*(U) + \b*(V) $) circle (3pt)
                node[xshift=4mm,yshift=-3mm] {$\a\vect u \ifnum\b<0 \b\vect v\else+\b\vect v\fi$};
            }

            \draw[->, very thick]
            (0,0,0) -- (U) node[above right] {$\vect u$};

            \draw[->, very thick]
            (0,0,0) -- (V) node[below left] {$\vect v$};

            \node at (0,4,7)
            {$\operatorname{span}\{\vect u,\vect v\}$};
        \end{tikzpicture}

        \subcaption{Udspændet af to ikke-parallelle vektorer i $\mathbb{R}^3$ er en plan gennem origo.}
    \end{minipage}

    \caption{To eksempler på den geometriske fortolkning af udspænd}
    \label{fig:span_geometry}
\end{figure}

\section{Matrixligninger}
I de foregående afsnit så vi på forskellige måder at beskrive lineære sammenhænge på. Vi så for eksempel på enkelte lineære ligninger som  
\[
2x_1 + 3x_2 = 5,
\]
og på lineære ligningssystemer som  
\[
\begin{aligned}
x_1 + 2x_2 &= 4 \\
3x_1 - x_2 &= 2.
\end{aligned}
\]
Vi skrev også disse systemer som vektorligninger, fx  
\[
x_1 \begin{bmatrix} 1 \\ 3 \end{bmatrix} + x_2 \begin{bmatrix} 2 \\ -1 \end{bmatrix} = \begin{bmatrix} 4 \\ 2 \end{bmatrix}.
\]

Nu er det tid til at introducere en ny og stærk form: \textbf{matrixligningen}. Denne kompakte skrivemåde gør det muligt at håndtere lineære ligningssystemer effektivt med matrixnotation. På denne form bliver vores system til  
\[
A \vect{x} = \vect{b},
\]
hvor \(A\) er en matrix, \(\vect{x}\) er en vektor af variable, og \(\vect{b}\) er resultatvektoren. Matrixligninger giver os en struktureret, algebraisk tilgang til at løse ligningssystemer.

\begin{definition}{Matrixligning}
  Hvis $A$ er en $m \times n$-matrix med søjlerne $\vect{a}_1, \ldots, \vect{a}_n$, og hvis $\vect{x} \in \mathbb{R}^n$, så er produktet af $A$ og $\vect{x}$, som skrives $A \vect{x}$, givet ved

\[
A \vect{x} = \left[\begin{array}{llll}
\vect{a}_1 & \vect{a}_2 & \ldots & \vect{a}_n
\end{array}\right] \left[\begin{array}{c}
x_1 \\
x_2 \\
\vdots \\
x_n
\end{array}\right] = x_1 \vect{a}_1 + x_2 \vect{a}_2 + \cdots + x_n \vect{a}_n
\]
    
\end{definition}

\begin{example}
    \[
\begin{aligned}
{\left[\begin{array}{lll}
4 & 1 & 2 \\
8 & 0 & 3
\end{array}\right]\left[\begin{array}{l}
2 \\
1 \\
4
\end{array}\right] } & =2\left[\begin{array}{l}
4 \\
8
\end{array}\right]+1\left[\begin{array}{l}
1 \\
0
\end{array}\right]+4\left[\begin{array}{l}
2 \\
3
\end{array}\right] \\
& =\left[\begin{array}{c}
8 \\
16
\end{array}\right]+\left[\begin{array}{l}
1 \\
0
\end{array}\right]+\left[\begin{array}{c}
8 \\
12
\end{array}\right] \\
& =\left[\begin{array}{l}
17 \\
28
\end{array}\right]
\end{aligned}
\]
\end{example}

Ved at bygge lidt videre på hovedresultatet fra det foregående afsnit får vi følgende sætning.

\begin{theorem}
    Hvis $A$ er en $m \times n$-matrix med søjlerne $\vect{a}_1, \ldots, \vect{a}_n \in \mathbb{R}^m$, og hvis $\vect{b} \in \mathbb{R}^m$, så har matrixligningen

\[
A \vect{x} = \vect{b}
\]

samme løsningsmængde som vektorligningen

\[
x_1 \vect{a}_1 + x_2 \vect{a}_2 + \cdots + x_n \vect{a}_n = \vect{b}
\]

som igen har samme løsningsmængde som det lineære ligningssystem med totalmatrixen

\[
\left[\begin{array}{lllll}
\vect{a}_1 & \vect{a}_2 & \ldots & \vect{a}_n & \vect{b}
\end{array}\right].
\]

\end{theorem}

Sætningen er vigtig, fordi den binder de tre måder at beskrive og løse lineære ligningssystemer på sammen: matrixligninger, vektorligninger og ligningssystemer med totalmatricer. Den viser, at uanset hvilken form vi bruger, fører de alle til samme løsningsmængde.

Når matrixligningen \( A \vect{x} = \vect{b} \) er ækvivalent med vektorligningen \( x_1 \vect{a}_1 + x_2 \vect{a}_2 + \cdots + x_n \vect{a}_n = \vect{b} \), kan vi tolke det at løse en matrixligning som at finde den rette kombination af søjlerne i \( A \), der giver \( \vect{b} \). Denne sammenhæng gør det muligt at anskueliggøre problemet ved hjælp af vektorrum og linearkombinationer.

Sætningen viser desuden, at den samme proces afspejles i totalmatrixen, hvor rækkereduktion afslører løsningen gennem elementære rækkeoperationer. Uanset om vi angriber problemet algebraisk, geometrisk eller algoritmisk, har vi altså med den samme underliggende struktur at gøre. Dette samlende perspektiv forenkler vores tilgang til at løse ligningssystemer og giver fleksibilitet i, hvordan vi vælger at beskrive og behandle problemet.

\begin{remark}
    Ligningen $A \vect{x} = \vect{b}$ har en løsning, hvis og kun hvis $\vect{b}$ er en linearkombination af søjlerne i $A$.
\end{remark}

Følgende sætning fortæller os, hvornår søjlevektorerne i en $m \times n$-matrix (søjlerne er vektorer i $\mathbb{R}^m$, da der er $m$ rækker) kan bruges til at frembringe hele $\mathbb{R}^m$. Den opsummerer i bund og grund flere ting, vi allerede har set i forskellige sammenhænge.

\begin{theorem}
Lad \( A \) være en \( m \times n \)-matrix. Så er følgende udsagn ækvivalente:
\begin{enumerate}[label=(\alph*)]
    \item For hver \( \vect{b} \in \mathbb{R}^m \) har ligningen \( A \vect{x} = \vect{b} \) en løsning.
    \item Hver \( \vect{b} \in \mathbb{R}^m \) er en linearkombination af søjlerne i \( A \).
    \item Søjlerne i \( A \) udspænder \( \mathbb{R}^m \).
    \item Matricen \( A \) har en pivotposition i hver række.
\end{enumerate}
\end{theorem}

Sætningen viser, at flere tilsyneladende forskellige idéer faktisk er ækvivalente. Om matrixligningen \( A \vect{x} = \vect{b} \) har en løsning (udsagn (a)), afhænger af, om søjlerne i \( A \) kan danne en hvilken som helst vektor i \( \mathbb{R}^m \) (udsagn (b)). Geometrisk betyder det, at søjlerne udspænder hele rummet \( \mathbb{R}^m \) (udsagn (c)).

Endelig giver tilstedeværelsen af en pivotposition i hver række (udsagn (d)) en algebraisk betingelse, der garanterer, at søjlernes udspænd dækker \( \mathbb{R}^m \), så \( A \vect{x} = \vect{b} \) altid har en løsning.

Denne ækvivalens er et stærkt værktøj, fordi den forbinder løsninger, linearkombinationer og geometriske fortolkninger og samtidig giver en konkret metode (at undersøge pivotpositionerne) til at efterprøve disse egenskaber.

\begin{custombox}{Række-vektor-reglen for beregning af \( A \vect{x} \)}
Forudsat at produktet \( A \vect{x} \) er defineret, er den \( i \)'te indgang i \( A \vect{x} \) summen af produkterne af de tilsvarende indgange fra række \( i \) i \( A \) og vektoren \( \vect{x} \). Med andre ord er den \( i \)'te indgang i \( A \vect{x} \) prikproduktet af den vektor, der udgør den \( i \)'te række i \( A \), og vektoren \( \vect{x} \).
    
\end{custombox}


Reglen forenkler matrixmultiplikation ved at dele den op i velkendte prikprodukter. Der findes også en komplementær tolkning via søjlerne: Hvis søjlerne i $A$ er $\vect{a}_1,\ldots,\vect{a}_n$ og $\vect{x}=(x_1,\ldots,x_n)^T$, så er
\[
A\vect{x}=x_1\vect{a}_1+\cdots+x_n\vect{a}_n.
\]
$A\vect{x}$ er altså en linearkombination af \emph{søjlerne} i $A$, mens hver enkelt indgang i produktet er et prikprodukt med en række i $A$.

\begin{example} Lad

\[
A=\left[\begin{array}{ccc}
1 & -1 & 4 \\
5 & 0 & 2
\end{array}\right], \text { og } \vect{x}=\left[\begin{array}{l}
2 \\
1 \\
5
\end{array}\right]
\]

Så er

\[
\begin{aligned}
{\left[\begin{array}{ccc}
1 & -1 & 4 \\
5 & 0 & 2
\end{array}\right]\left[\begin{array}{l}
2 \\
1 \\
5
\end{array}\right] } & =\left[\begin{array}{c}
1 \cdot 2+(-1) \cdot 1+4 \cdot 5 \\
5 \cdot 2+0 \cdot 1+2 \cdot 5
\end{array}\right] \\
& =\left[\begin{array}{c}
21 \\
20
\end{array}\right]
\end{aligned}
\]
\end{example}

Bemærk, at antallet af søjler i $A$ skal være lig med antallet af rækker i $\vect{x}$, da prikproduktet ellers ikke ville give mening!

Den næste sætning indfanger to vigtige egenskaber ved multiplikation af en matrix med en vektor: \textbf{distributivitet} og \textbf{multiplikation med skalar}. Egenskaberne svarer til de velkendte regneregler i algebra, men gælder nu for matricer og vektorer.

\begin{theorem}
Hvis \( A \) er en \( m \times n \)-matrix, \( \vect{u} \) og \( \vect{v} \) er vektorer i \( \mathbb{R}^n \), og \( c \) er en skalar, så gælder:

\begin{enumerate}[label=(\alph*)]
    \item \( A(\vect{u} + \vect{v}) = A \vect{u} + A \vect{v} \)  


    \item \( A(c \vect{u}) = c(A \vect{u}) \)  

\end{enumerate}
\end{theorem}

Egenskab (a) viser, at matrixmultiplikation er distributiv over vektoraddition. Det betyder, at det at gange \( A \) med summen af to vektorer giver det samme som at gange \( A \) med hver vektor for sig og derefter lægge resultaterne sammen.     Egenskab (b) viser, at multiplikation med skalar kommuterer med matrixmultiplikation. At gange en vektor med en skalar først og derefter anvende matricen giver det samme resultat som først at anvende matricen og derefter gange den resulterende vektor med skalaren.

\begin{example} Lad

\[
A = \left[\begin{array}{cccc}
1 & 5 & -2 & 0 \\
-3 & 1 & 9 & -5 \\
4 & -8 & -1 & 7
\end{array}\right], \quad \vect{p} = \left[\begin{array}{c}
3 \\
-2 \\
0 \\
-4
\end{array}\right], \quad \vect{b} = \left[\begin{array}{c}
-7 \\
9 \\
0
\end{array}\right]
\]

Det kan vises, at $\vect{p}$ er en løsning til $A \vect{x} = \vect{b}$. Brug dette til at skrive $\vect{b}$ som en linearkombination af søjlerne i $A$.
    
\end{example}  

\begin{solution} Da $\vect{p}$ er en løsning til $A \vect{x} = \vect{b}$, har vi:
    \[
    A \vect{p} = \vect{b}.
    \]
    
    Vi kan udtrykke $\vect{b}$ som en linearkombination af søjlerne i $A$ ved hjælp af indgangene i $\vect{p}$. Lad $\vect{a}_1$, $\vect{a}_2$, $\vect{a}_3$ og $\vect{a}_4$ betegne søjlerne i $A$, så:
    \[
    A = \left[ \vect{a}_1 \ \big| \ \vect{a}_2 \ \big| \ \vect{a}_3 \ \big| \ \vect{a}_4 \right].
    \]
    
    Dermed er
    \[
    \vect{b} = p_1 \vect{a}_1 + p_2 \vect{a}_2 + p_3 \vect{a}_3 + p_4 \vect{a}_4.
    \]
    
    Nu beregner vi hvert led for sig.
    
    \textbf{Beregn $p_1 \vect{a}_1$:}
    \[
    p_1 \vect{a}_1 = 3 \left[\begin{array}{c}
    1 \\
    -3 \\
    4
    \end{array}\right] = \left[\begin{array}{c}
    3 \\
    -9 \\
    12
    \end{array}\right].
    \]
    
    \textbf{Beregn $p_2 \vect{a}_2$:}
    \[
    p_2 \vect{a}_2 = (-2) \left[\begin{array}{c}
    5 \\
    1 \\
    -8
    \end{array}\right] = \left[\begin{array}{c}
    -10 \\
    -2 \\
    16
    \end{array}\right].
    \]
    
    \textbf{Beregn $p_3 \vect{a}_3$:}
    \[
    p_3 \vect{a}_3 = 0 \left[\begin{array}{c}
    -2 \\
    9 \\
    -1
    \end{array}\right] = \left[\begin{array}{c}
    0 \\
    0 \\
    0
    \end{array}\right].
    \]
    
    \textbf{Beregn $p_4 \vect{a}_4$:}
    \[
    p_4 \vect{a}_4 = (-4) \left[\begin{array}{c}
    0 \\
    -5 \\
    7
    \end{array}\right] = \left[\begin{array}{c}
    0 \\
    20 \\
    -28
    \end{array}\right].
    \]
    
    \textbf{Læg de beregnede vektorer sammen for at finde $\vect{b}$:}
    \[
    \vect{b} = p_1 \vect{a}_1 + p_2 \vect{a}_2 + p_3 \vect{a}_3 + p_4 \vect{a}_4 = \left[\begin{array}{c}
    3 \\
    -9 \\
    12
    \end{array}\right] + \left[\begin{array}{c}
    -10 \\
    -2 \\
    16
    \end{array}\right] + \left[\begin{array}{c}
    0 \\
    0 \\
    0
    \end{array}\right] + \left[\begin{array}{c}
    0 \\
    20 \\
    -28
    \end{array}\right].
    \]
    
    Nu beregner vi summen trin for trin.
    
    \textbf{Trin 1: Læg $p_1 \vect{a}_1$ og $p_2 \vect{a}_2$ sammen:}
    \[
    \left[\begin{array}{c}
    3 \\
    -9 \\
    12
    \end{array}\right] + \left[\begin{array}{c}
    -10 \\
    -2 \\
    16
    \end{array}\right] = \left[\begin{array}{c}
    -7 \\
    -11 \\
    28
    \end{array}\right].
    \]
    
    \textbf{Trin 2: Læg $p_3 \vect{a}_3$ (som er nul) til resultatet:}
    \[
    \left[\begin{array}{c}
    -7 \\
    -11 \\
    28
    \end{array}\right] + \left[\begin{array}{c}
    0 \\
    0 \\
    0
    \end{array}\right] = \left[\begin{array}{c}
    -7 \\
    -11 \\
    28
    \end{array}\right].
    \]
    
    \textbf{Trin 3: Læg $p_4 \vect{a}_4$ til resultatet:}
    \[
    \left[\begin{array}{c}
    -7 \\
    -11 \\
    28
    \end{array}\right] + \left[\begin{array}{c}
    0 \\
    20 \\
    -28
    \end{array}\right] = \left[\begin{array}{c}
    -7 \\
    9 \\
    0
    \end{array}\right].
    \]
    
    \textbf{Endeligt resultat:}
    \[
    \vect{b} = \left[\begin{array}{c}
    -7 \\
    9 \\
    0
    \end{array}\right].
    \]
    
    \textbf{Konklusion:}
    
    Vi har udtrykt $\vect{b}$ som en linearkombination af søjlerne i $A$ ved hjælp af indgangene i $\vect{p}$:
    \[
    \vect{b} = 3 \vect{a}_1 - 2 \vect{a}_2 + 0 \vect{a}_3 - 4 \vect{a}_4.
    \]
\end{solution}

\section{Løsningsmængder for lineære ligningssystemer}
Løsningsmængder for lineære ligningssystemer spiller en central rolle i lineær algebra og dukker op igen i mange sammenhænge i faget. I dette afsnit bruger vi vektornotation til at give klare, eksplicitte og geometriske beskrivelser af disse løsningsmængder.

Et lineært ligningssystem kaldes \textit{homogent}, hvis det kan skrives på formen

\[
A \vect{x} = \vect{0}
\]

hvor $A$ er en $m \times n$-matrix, og $\vect{0}$ er nulvektoren i $\mathbb{R}^m$. Et sådant system har altid den \textit{trivielle løsning}, nemlig $\vect{x} = \vect{0}$ (nulvektoren i $\mathbb{R}^n$, ikke $\mathbb{R}^m$).

Da homogene systemer altid har den trivielle løsning, er det interessante spørgsmål, om de har \textit{ikke-trivielle løsninger}.

\begin{remark}
    Den homogene ligning $A \vect{x} = \vect{0}$ har en ikke-triviel løsning, hvis og kun hvis ligningen har mindst én fri variabel.
\end{remark}

\begin{example}
    Afgør, om følgende homogene system har en ikke-triviel løsning. Hvis det har, så beskriv løsningsmængden ved hjælp af den eller de frie variable.
\[
\begin{aligned}
2 x_1+x_2-3 x_3 & =0 \\
x_1-x_2+x_3 & =0 \\
-2 x_1+5 x_2-7 x_3 & =0
\end{aligned}
\]

\end{example}

\begin{solution}
    Vi lader $A$ være systemets koefficientmatrix og reducerer totalmatrixen [ $\left.\begin{array}{cc}A & 0\end{array}\right]$ til echelonform (bemærk, at vi sætter den anden ligning først for at gøre beregningen lidt enklere):

\[
\begin{aligned}
{\left[\begin{array}{cccc}
1 & -1 & 1 & 0 \\
2 & 1 & -3 & 0 \\
-2 & 5 & -7 & 0
\end{array}\right] } & \sim\left[\begin{array}{cccc}
1 & -1 & 1 & 0 \\
0 & 3 & -5 & 0 \\
0 & 3 & -5 & 0
\end{array}\right] \\
& \sim\left[\begin{array}{cccc}
1 & -1 & 1 & 0 \\
0 & 3 & -5 & 0 \\
0 & 0 & 0 & 0
\end{array}\right]
\end{aligned}
\]

så vi ser, at $x_3$ er en fri variabel, og at der derfor findes ikke-trivielle løsninger. Nu bestemmer vi den reducerede echelonform:

\[
\begin{aligned}
{\left[\begin{array}{cccc}
1 & -1 & 1 & 0 \\
0 & 3 & -5 & 0 \\
0 & 0 & 0 & 0
\end{array}\right] } & \sim\left[\begin{array}{cccc}
1 & -1 & 1 & 0 \\
0 & 1 & -\frac{5}{3} & 0 \\
0 & 0 & 0 & 0
\end{array}\right] \\
& \sim\left[\begin{array}{cccc}
1 & 0 & -\frac{2}{3} & 0 \\
0 & 1 & -\frac{5}{3} & 0 \\
0 & 0 & 0 & 0
\end{array}\right]
\end{aligned}
\]

Det giver os ligningerne

\[
\begin{aligned}
& x_1-\frac{2}{3} x_3=0 \quad \Longrightarrow x_1=\frac{2}{3} x_3 \\
& x_2-\frac{5}{3} x_3=0 \quad \Longrightarrow x_2=\frac{5}{3} x_3
\end{aligned}
\]


På vektorform skrives løsningen $\vect{x}$ til ligningen $A \vect{x}=0$ som

\[
\vect{x}=\left[\begin{array}{l}
x_1 \\
x_2 \\
x_3
\end{array}\right]=\left[\begin{array}{c}
\frac{2}{3} x_3 \\
\frac{5}{3} x_3 \\
x_3
\end{array}\right]=x_3\left[\begin{array}{c}
\frac{2}{3} \\
\frac{5}{3} \\
1
\end{array}\right]
\]


Da $x_3$ kan være hvad som helst, beskriver denne løsningsmængde geometrisk den linje i $\mathbb{R}^3$, der går fra origo gennem punktet $\left(\frac{2}{3}, \frac{5}{3}, 1\right)$.
\end{solution}


\begin{custombox}{Parameterfremstilling på vektorform}
    Når en løsningsmængde er beskrevet eksplicit med vektorer (som i det foregående eksempel), siger vi, at løsningen er givet som en \textbf{parameterfremstilling på vektorform}.
\end{custombox}

\begin{example} Beskriv alle løsninger til den homogene ligning


\[
    x_1-2 x_2-5 x_3=0
\]

\end{example}

\begin{solution}

Hvis vi skriver systemet som en matrix, ser vi, at $x_2$ og $x_3$ er frie variable, og at $x_1$ er en basisvariabel med $x_1=2 x_2+5 x_3$. Den fuldstændige løsning er derfor

$$
\vect{x}=\left[\begin{array}{l}
x_1 \\
x_2 \\
x_3
\end{array}\right]=\left[\begin{array}{c}
2 x_2+5 x_3 \\
x_2 \\
x_3
\end{array}\right]=x_2\left[\begin{array}{l}
2 \\
1 \\
0
\end{array}\right]+x_3\left[\begin{array}{l}
5 \\
0 \\
1
\end{array}\right]
$$

hvor den yderste højre del af ligningen er løsningsmængdens \textbf{parameterfremstilling på vektorform}. I dette tilfælde er hele løsningsmængden en plan i $\mathbb{R}^3$, som indeholder de to linjer fra $(0,0,0)$ til $(2,1,0)$ og fra $(0,0,0)$ til $(5,0,1)$.
\end{solution}

Når et \textit{inhomogent} lineært ligningssystem har mange løsninger, kan den fuldstændige løsning skrives som en parameterfremstilling på vektorform, der består af én vektor plus vilkårlige linearkombinationer af vektorer, som opfylder det tilhørende homogene system.

\begin{example} Vi ser igen på det homogene system fra begyndelsen af dette afsnit, men denne gang er det inhomogent (dvs. højresiden består ikke udelukkende af nuller):

\[
\begin{aligned}
2 x_1 + x_2 - 3 x_3 & = 2 \\
x_1 - x_2 + x_3 & = 1 \\
-2 x_1 + 5 x_2 - 7 x_3 & = -2
\end{aligned}
\]

Som før udfører vi rækkeoperationer på totalmatrixen \(\left[\begin{array}{cc} A & \vect{b} \end{array}\right]\), hvor \( A \) er den samme koefficientmatrix og

\[
\vect{b} = \begin{bmatrix}
2 \\
1 \\
-2
\end{bmatrix}
\]

og vi får, at

\[
\left[\begin{array}{cccc}
1 & -1 & 1 & 1 \\
2 & 1 & -3 & 2 \\
-2 & 5 & -7 & -2
\end{array}\right]
\sim
\left[\begin{array}{cccc}
1 & 0 & -\frac{2}{3} & 1 \\
0 & 1 & -\frac{5}{3} & 0 \\
0 & 0 & 0 & 0
\end{array}\right]
\]

Dermed er

\[
\begin{aligned}
x_1 &= 1 + \frac{2}{3} x_3, \text{ og } \\
x_2 &= \frac{5}{3} x_3
\end{aligned}
\]

Dette kan skrives som

\[
\vect{x} = \begin{bmatrix}
x_1 \\
x_2 \\
x_3
\end{bmatrix} =
\begin{bmatrix}
1 + \frac{2}{3} x_3 \\
\frac{5}{3} x_3 \\
x_3
\end{bmatrix}
= \begin{bmatrix}
1 \\
0 \\
0
\end{bmatrix} + x_3 \begin{bmatrix}
\frac{2}{3} \\
\frac{5}{3} \\
1
\end{bmatrix}
\]

som har den form, vi påstod. Det er ligningen for linjen gennem punktet \((1,0,0)\), der er parallel med linjen fra \((0,0,0)\) gennem \(\left(\frac{2}{3}, \frac{5}{3}, 1\right)\).

\end{example}

Vi opsummerer den generelle situation med følgende sætning. Konklusionen kan måske bedst opfattes som en forskydning af alle de homogene løsninger med vektoren $\vect{p}$.

\begin{theorem}
    Antag, at $A \vect{x} = \vect{b}$ er konsistent for et givet $\vect{b}$, og lad $\vect{p}$ være en løsning. Så er løsningsmængden for $A \vect{x} = \vect{b}$ mængden af alle vektorer på formen
    \[
    \vect{w} = \vect{p} + \vect{v}_h,
    \]
    hvor $\vect{v}_h$ er en vilkårlig løsning til den homogene ligning $A \vect{x} = \vect{0}$.
\end{theorem}

\begin{example} Skriv den fuldstændige løsning til

\[
x_1-2 x_2-5 x_3=3
\]

som en parameterfremstilling på vektorform. Hvordan ser denne løsningsmængde geometrisk ud sammenlignet med løsningsmængden for ligningen $x_1-2 x_2-5 x_3=0$, som vi så tidligere?
\end{example}

\begin{solution}
        Vi løser ligningen:
        \[
        x_1 - 2 x_2 - 5 x_3 = 3
        \]
        
        For at udtrykke den fuldstændige løsning løser vi med hensyn til \( x_1 \) udtrykt ved \( x_2 \) og \( x_3 \):
    
        \[
        \begin{aligned}
        x_1 &= 3 + 2x_2 + 5x_3, \\
        x_2 &= x_2, \\
        x_3 &= x_3.
        \end{aligned}
        \]
        
        Vi kan skrive løsningsmængden som en \textit{parameterfremstilling på vektorform}:
        \[
        \begin{bmatrix}
        x_1 \\
        x_2 \\
        x_3
        \end{bmatrix}
        =
        \begin{bmatrix}
        3 \\
        0 \\
        0
        \end{bmatrix}
        +
        x_2
        \begin{bmatrix}
        2 \\
        1 \\
        0
        \end{bmatrix}
        +
        x_3
        \begin{bmatrix}
        5 \\
        0 \\
        1
        \end{bmatrix}, \quad \text{for alle } x_2, x_3 \in \mathbb{R}.
        \]
        
        \textbf{Geometrisk fortolkning:} Løsningsmængden er en plan i \(\mathbb{R}^3\). Planen er parallel med den plan, der er bestemt af den homogene ligning:
        \[
        x_1 - 2 x_2 - 5 x_3 = 0
        \]
        
        Den homogene løsningsmængde kan skrives som:
        \[
        \begin{bmatrix}
        x_1 \\
        x_2 \\
        x_3
        \end{bmatrix}
        =
        x_2
        \begin{bmatrix}
        2 \\
        1 \\
        0
        \end{bmatrix}
        +
        x_3
        \begin{bmatrix}
        5 \\
        0 \\
        1
        \end{bmatrix}, \quad \text{for alle } x_2, x_3 \in \mathbb{R}.
        \]
        
        Når vi sammenligner de to løsningsmængder, ser vi, at den oprindelige løsning er den homogene løsning forskudt med vektoren \(\begin{bmatrix} 3 \\ 0 \\ 0 \end{bmatrix}\). Geometrisk betyder det, at planen bestemt af \( x_1 - 2 x_2 - 5 x_3 = 3 \) er parallel med, men ikke den samme som, planen bestemt af \( x_1 - 2 x_2 - 5 x_3 = 0 \). De to planer er forskudt 3 enheder i forhold til hinanden langs \( x_1 \)-aksen. Den fuldstændige løsning til ligningen \( x_1 - 2 x_2 - 5 x_3 = 3 \) er en plan i \(\mathbb{R}^3\), der er parallel med løsningsmængden for den homogene ligning \( x_1 - 2 x_2 - 5 x_3 = 0 \), men forskudt væk fra origo. \end{solution}

        \begin{example} Find en parameterfremstilling for linjen gennem $\vect{a}=\left[\begin{array}{c}3 \\ -2\end{array}\right]$, der også er parallel med $\vect{b}=\left[\begin{array}{c}-7 \\ 6\end{array}\right]$.
        \end{example}

        \begin{solution}
            Linjen gennem $\vect{a}$ parallel med $\vect{b}$ er givet ved ligningen $\vect{r}=\vect{a}+t \vect{b}$, hvor $t$ er en parameter. Ved at indsætte de givne vektorer får vi
            \[
            \begin{aligned}
            \vect{r} &=\left[\begin{array}{c}
            3 \\
            -2
            \end{array}\right]+t\left[\begin{array}{c}
            -7 \\
            6
            \end{array}\right] \\
            &=\left[\begin{array}{c}
            3-7 t \\
            -2+6 t
            \end{array}\right]
            \end{aligned}
            \]
            Dette er en parameterfremstilling for linjen gennem $\vect{a}$, der er parallel med $\vect{b}$.
        \end{solution}

\section{Lineær uafhængighed}
Vi afslutter dette afsnit med at introducere et af de vigtigste begreber i lineær algebra: \textit{lineær uafhængighed}. Begrebet er grundlæggende for at forstå vektorers struktur og deres indbyrdes sammenhænge i vektorrum.

\begin{definition}
    En indekseret mængde af vektorer $\left\{\vect{v}_1, \ldots, \vect{v}_p\right\} \subset \mathbb{R}^n$ kaldes \textbf{lineært uafhængig}, hvis vektorligningen
    
    \[
    x_1 \vect{v}_1 + \cdots + x_p \vect{v}_p = \vect{0}
    \]
    
    kun har den trivielle løsning, dvs. hvis den eneste løsning er $\left(x_1, \ldots, x_p\right) = (0, \ldots, 0)$. Tilsvarende kaldes mængden $\left\{\vect{v}_1, \ldots, \vect{v}_p\right\}$ \textbf{lineært afhængig}, hvis der findes vægte $c_1, \ldots, c_p$, som ikke alle er nul, så
    
    \[
    c_1 \vect{v}_1 + \cdots + c_p \vect{v}_p = \vect{0}
    \]
    
    En sådan ligning kalder vi en \textit{lineær afhængighedsrelation}, når vægtene ikke alle er nul.
\end{definition}
    
    \begin{remark}
        En mængde af vektorer kan ikke være både lineært uafhængig og lineært afhængig, men den skal være det ene af de to!
    \end{remark}

    At afgøre, om en mængde af vektorer er lineært uafhængig, svarer til at løse matrixligningen

    \[
    A \vect{x}=0,
    \]

hvor søjlerne i $A$ er de givne vektorer. Mængden af vektorer er lineært uafhængig, hvis og kun hvis den eneste løsning er $\vect{x}=\vect{0}$. Ellers findes der en lineær afhængighedsrelation.

    \begin{example}
        Afgør, om følgende vektorer i $\mathbb{R}^3$ er lineært uafhængige:

        \[
        \vect{v}_1 = \left[\begin{array}{l}
        2 \\
        1 \\
        2
        \end{array}\right], \quad \vect{v}_2 = \left[\begin{array}{l}
        1 \\
        1 \\
        1
        \end{array}\right], \quad \vect{v}_3 = \left[\begin{array}{l}
        0 \\
        1 \\
        1
        \end{array}\right]
        \]
    \end{example}

    \begin{solution}
        Vi begynder med at reducere den tilhørende totalmatrix til echelonform:

\[
\begin{aligned}
{\left[\begin{array}{llll}
2 & 1 & 0 & 0 \\
1 & 1 & 1 & 0 \\
2 & 1 & 1 & 0
\end{array}\right] } & \sim\left[\begin{array}{llll}
1 & 1 & 1 & 0 \\
2 & 1 & 0 & 0 \\
2 & 1 & 1 & 0
\end{array}\right] \\
& \sim\left[\begin{array}{cccc}
1 & 1 & 1 & 0 \\
0 & -1 & -2 & 0 \\
0 & -1 & -1 & 0
\end{array}\right] \\
& \sim\left[\begin{array}{cccc}
1 & 1 & 1 & 0 \\
0 & -1 & -2 & 0 \\
0 & 0 & 1 & 0
\end{array}\right]
\end{aligned}
\]

hvoraf vi ser, at der findes en løsning. Nu vil vi vide, om løsningen er den trivielle løsning $(0,0,0)$. Normalt ville vi fortsætte med rækkeoperationer, indtil vi når den reducerede echelonform, men vi kan gøre det smartere. Bemærk først og fremmest, at der ikke er nogen »dårlige« rækker (så der findes mindst én løsning), hvilket vi også forventer, da et homogent system altid har mindst den trivielle løsning. Da der ikke er nogen frie variable, ser vi, at der er præcis én løsning. Det fortæller os, at løsningen må være

\[
\vect{x}=\left[\begin{array}{l}
x_1 \\
x_2 \\
x_3
\end{array}\right]=\left[\begin{array}{l}
0 \\
0 \\
0
\end{array}\right]
\]

Mængden af vektorer $\left\{\vect{v}_1, \vect{v}_2, \vect{v}_3\right\}$ er altså lineært uafhængig.
\end{solution}

Generelt er det meget nyttigt at kunne afgøre, om en mængde af vektorer er lineært uafhængig, så det er godt at have
nogle sætninger, der hurtigt klarer problemet i visse specialtilfælde.

\begin{theorem}
    \begin{enumerate}[(a)]
        \item Hvis en mængde af vektorer indeholder nulvektoren, er mængden lineært afhængig.
        \item Hvis en mængde af vektorer indeholder et skalarmultiplum af en anden vektor i mængden, er mængden lineært afhængig.
        \item Hvis en mængde af vektorer indeholder flere vektorer, end der er indgange i hver vektor, er mængden lineært afhængig.
    \end{enumerate}
\end{theorem}

Sætningen giver hurtige test for lineær afhængighed. Mere generelt er en mængde lineært afhængig, hvis en af dens vektorer er en linearkombination af de andre. Vi kan formalisere denne idé i følgende korollar.

\begin{corollary}
    En indekseret mængde $S = \left\{\vect{v}_1, \ldots, \vect{v}_p\right\}$ af to eller flere vektorer er lineært afhængig, hvis og kun hvis mindst én af vektorerne i $S$ er en linearkombination af de andre. Hvis $S$ er lineært afhængig, og $\vect{v}_1 \neq \vect{0}$, så er en af vektorerne $\vect{v}_j$ (med $1 < j \leq p$) en linearkombination af de foregående vektorer $\vect{v}_1, \ldots, \vect{v}_{j-1}$.
\end{corollary}

\begin{remark}
    Korollaret fortæller os, at hvis en mængde af vektorer er lineært afhængig, er mindst én af vektorerne overflødig, da den kan udtrykkes som en linearkombination af de andre. Det er denne overflødighed, der giver den lineære afhængighed. Det betyder ikke, at hver vektor i $S$ er en linearkombination af de andre, men kun at \textit{mindst én} vektor i en lineært afhængig mængde er en linearkombination af andre.
\end{remark}

\begin{example} Lad
    \[
    \vect{u} = \left[\begin{array}{l}
    2 \\
    1 \\
    1
    \end{array}\right], \quad \vect{v} = \left[\begin{array}{l}
    4 \\
    3 \\
    5
    \end{array}\right], \quad \vect{w} = \left[\begin{array}{l}
    1 \\
    1 \\
    1
    \end{array}\right], \quad \vect{z} = \left[\begin{array}{l}
    0 \\
    0 \\
    1
    \end{array}\right].
    \]

    \begin{enumerate}[(a)]
        \item Er nogen af parrene af disse vektorer (fx $\{\vect{u}, \vect{v}\}$) lineært afhængige? Begrund dit svar.
        \item Fortæller svaret på del (a) os, at $\{\vect{u}, \vect{v}, \vect{w}, \vect{z}\}$ er lineært uafhængig?
        \item Er $\{\vect{u}, \vect{v}, \vect{w}, \vect{z}\}$ lineært afhængig? Du bør kunne besvare spørgsmålet uden at regne.
    \end{enumerate}
\end{example}

\begin{solution}
    \begin{enumerate}[(a)]
        \item 
        For at afgøre, om nogen af parrene af disse vektorer er lineært afhængige, undersøger vi, om den ene vektor er et skalarmultiplum af den anden.

        \textbf{Undersøgelse af parret $\{\vect{u}, \vect{v}\}$:} Antag, at der findes en skalar $k$, så:
    
        \[
        \vect{v} = k \vect{u}.
        \]
        
        Vi beregner $k$ ud fra den første komposant:
        \[
        k = \dfrac{v_1}{u_1} = \dfrac{4}{2} = 2.
        \]
        
        Vi kontrollerer med den anden komposant:
        \[
        v_2 = k u_2 \implies 3 = 2 \times 1 \implies 3 = 2.
        \]
        
        Da $3 \ne 2$, er $\vect{v}$ ikke et skalarmultiplum af $\vect{u}$.
        
        \textbf{Undersøgelse af parret $\{\vect{u}, \vect{w}\}$:} Antag, at der findes en skalar $k$, så:
        \[
        \vect{u} = k \vect{w}.
        \]
        
        Vi beregner $k$ ud fra den første komposant:
        \[
        k = \dfrac{u_1}{w_1} = \dfrac{2}{1} = 2.
        \]
        
        Vi kontrollerer med den anden komposant:
        \[
        u_2 = k w_2 \implies 1 = 2 \times 1 \implies 1 = 2.
        \]
        
        Da $1 \ne 2$, er $\vect{u}$ ikke et skalarmultiplum af $\vect{w}$.
        
        \textbf{Undersøgelse af parret $\{\vect{w}, \vect{z}\}$:} Da $z_1 = 0$ og $w_1 = 1$, fører antagelsen $\vect{w} = k \vect{z}$ til:
        \[
        w_1 = k z_1 \implies 1 = k \times 0 \implies 1 = 0,
        \] hvilket er en modstrid.
    
        \item Nej, at intet par er lineært afhængigt, betyder ikke, at hele mængden er lineært uafhængig. Lineær uafhængighed kræver, at den eneste løsning til:
        \[
        c_1 \vect{u} + c_2 \vect{v} + c_3 \vect{w} + c_4 \vect{z} = \vect{0}
        \]
        er $c_1 = c_2 = c_3 = c_4 = 0$.
        
        \item Ja, mængden er lineært afhængig. I $\mathbb{R}^3$ er enhver mængde med mere end tre vektorer nødvendigvis lineært afhængig, fordi det største antal lineært uafhængige vektorer i $\mathbb{R}^3$ er tre.
    \end{enumerate}
        
        \textbf{Konklusion:}
        
        \begin{enumerate}[(a)]
            \item Intet par af disse vektorer er lineært afhængigt.
            \item Svaret på del (a) garanterer ikke, at hele mængden er lineært uafhængig.
            \item Mængden $\{\vect{u}, \vect{v}, \vect{w}, \vect{z}\}$ er lineært afhængig, uden at der er brug for yderligere beregninger.
        \end{enumerate}

\end{solution}
